% !TeX root = ../main.tex
\chapter{THỰC NGHIỆM VÀ KẾT QUẢ}
\label{chap:thucnghiem}

\section{Thiết lập thực nghiệm}
\label{sec:thietlap}

\subsection{Môi trường phần cứng và phần mềm}

\begin{table}[H]
\centering
\caption{Môi trường thực nghiệm}
\label{tab:moitruong}
\begin{tabular}{L{4.4cm} L{8.6cm}}
\toprule
\textbf{Thành phần} & \textbf{Cấu hình} \\
\midrule
Bộ xử lý đồ hoạ & NVIDIA Quadro M1000M, 2,15 GB VRAM \\
Hệ điều hành    & Windows 10 Pro \\
Python          & 3.13.12 \\
PyTorch         & 2.6.0+cu124~\cite{paszke2019pytorch} \\
torchvision     & 0.21.0+cu124 \\
CUDA            & 12.4 \\
Thư viện phân đoạn & segmentation-models-pytorch~\cite{iakubovskii2019smp} \\
Thư viện tăng cường & Albumentations~\cite{buslaev2020albumentations} \\
Thư viện xử lý ảnh & OpenCV~\cite{bradski2000opencv} \\
\bottomrule
\end{tabular}
\end{table}

Dung lượng bộ nhớ đồ hoạ chỉ 2,15 GB là ràng buộc chi phối nhiều quyết định thực
nghiệm và cần được nêu rõ để người đọc hiểu đúng bối cảnh. Ràng buộc này giải
thích trực tiếp việc chọn kiến trúc quy mô vừa (EfficientNet-B0 thay vì các biến
thể lớn hơn), kích thước lô nhỏ, và --- như phân tích ở \cref{subsec:confound}
--- cả một số khác biệt cấu hình giữa các phiên bản vốn không nằm trong ý đồ
thiết kế ban đầu.

\subsection{Chỉ số đánh giá và quy ước báo cáo}

Bài toán phân loại được đánh giá bằng \en{accuracy}, \en{precision}, \en{recall}
và $F_1$ theo \cref{eq:prf}, báo cáo cả ở mức tổng hợp lẫn mức từng lớp. Bài toán
phân đoạn dùng IoU và Dice theo \cref{eq:iou,eq:dicecoef}, kèm \en{precision} và
\en{recall} tính trên tập điểm ảnh.

Mọi chỉ số phân đoạn được tính \emph{theo từng ảnh rồi lấy trung bình} trên toàn
tập, chứ không gộp toàn bộ điểm ảnh của cả tập rồi tính một lần. Cách thứ nhất
cho mỗi ảnh trọng số như nhau, phù hợp khi diện tích tổn thương biến thiên mạnh
giữa các ảnh; cách thứ hai sẽ khiến các ảnh có vùng bệnh lớn chi phối kết quả.

Toàn bộ chỉ số báo cáo trong chương này được lấy trực tiếp từ ô kết quả của các
\en{notebook} thực nghiệm hoặc từ tệp kết quả đã lưu, và có thể truy vết về đúng
nguồn sinh ra chúng.

\section{Kết quả phân loại}
\label{sec:ketqua_cls}

\subsection{Diễn biến huấn luyện}

Mô hình được huấn luyện tối đa 50 epoch với cơ chế dừng sớm. Thực tế quá trình
dừng ở epoch 16 khi \texttt{val\_acc} không cải thiện trong 5 epoch liên tiếp.

\begin{figure}[H]
\centering
\includegraphics[width=\textwidth]{figures/cls_training_curves.png}
\caption[Đường cong huấn luyện mô hình phân loại]{Diễn biến hàm mất mát (trái) và
độ chính xác (phải) qua 16 epoch. Đường nét đứt đánh dấu độ chính xác kiểm định
tốt nhất 98,19\% đạt tại epoch 11.}
\label{fig:cls_curves}
\end{figure}

\begin{table}[H]
\centering
\caption{Diễn biến huấn luyện mô hình phân loại qua 16 epoch}
\label{tab:cls_history}
\small
\begin{tabular}{c r r r r l}
\toprule
\textbf{Epoch} & \textbf{train\_loss} & \textbf{train\_acc} & \textbf{val\_loss} & \textbf{val\_acc} & \textbf{Ghi chú} \\
\midrule
1  & 1,0963 & 0,5892 & 0,4712 & 0,8623 & \\
2  & 0,4369 & 0,8547 & 0,2316 & 0,9187 & \\
3  & 0,2740 & 0,9056 & 0,1820 & 0,9368 & \\
4  & 0,2293 & 0,9249 & 0,1222 & 0,9616 & \\
5  & 0,1906 & 0,9327 & 0,1175 & 0,9639 & \\
6  & 0,1602 & 0,9478 & 0,1017 & 0,9684 & \\
7  & 0,1487 & 0,9507 & 0,1356 & 0,9503 & lần đầu không cải thiện \\
8  & 0,1220 & 0,9597 & 0,0926 & 0,9729 & \\
9  & 0,1130 & 0,9613 & 0,0887 & 0,9729 & \\
10 & 0,0989 & 0,9662 & 0,0890 & 0,9752 & \\
\textbf{11} & 0,0977 & 0,9671 & 0,0811 & \textbf{0,9819} & \textbf{điểm kiểm tra tốt nhất} \\
12 & 0,0820 & 0,9729 & 0,0910 & 0,9774 & \\
13 & 0,0791 & 0,9771 & 0,1023 & 0,9684 & \\
14 & 0,0747 & 0,9758 & 0,1108 & 0,9729 & \\
15 & 0,0559 & 0,9800 & 0,1050 & 0,9616 & \\
16 & 0,0476 & 0,9858 & 0,0831 & 0,9684 & dừng sớm kích hoạt \\
\bottomrule
\end{tabular}
\end{table}

Diễn biến chia thành ba pha rõ rệt. \emph{Pha hội tụ nhanh} (epoch 1--6): hàm mất
mát kiểm định giảm từ 0,4712 xuống 0,1017, độ chính xác tăng vọt từ 86,23\% lên
96,84\% --- hiệu quả điển hình của học chuyển giao, mô hình chỉ cần thích ứng đặc
trưng sẵn có với miền dữ liệu mới. \emph{Pha tinh chỉnh} (epoch 7--11): cải thiện
chậm dần, đạt đỉnh 98,19\% tại epoch 11. \emph{Pha quá khớp} (epoch 12--16): dấu
hiệu rất rõ --- \texttt{train\_acc} tiếp tục tăng tới 98,58\% trong khi
\texttt{val\_acc} dao động rồi giảm về 96,84\%, đồng thời \texttt{val\_loss} tăng
từ 0,0811 lên 0,1108 tại epoch 14. Dừng sớm đã can thiệp đúng thời điểm.

Đáng chú ý là ở epoch 16, \texttt{train\_acc} (98,58\%) đã vượt \texttt{val\_acc}
(96,84\%) --- khoảng cách 1,74 điểm phần trăm. Khoảng cách này còn nhỏ, cho thấy
mức quá khớp ở giai đoạn cuối là nhẹ và không phá huỷ khả năng khái quát.

\subsection{Hiệu năng trên tập kiểm thử}

\begin{table}[H]
\centering
\caption{Hiệu năng phân loại trên tập kiểm thử (890 ảnh)}
\label{tab:cls_test}
\begin{tabular}{l c c}
\toprule
\textbf{Chỉ số} & \textbf{Kiểm định (tốt nhất)} & \textbf{Kiểm thử} \\
\midrule
Accuracy            & \textbf{98,19\%} & \textbf{97,19\%} \\
Precision (weighted) & ---             & 97,22\% \\
Recall (weighted)    & ---             & 97,19\% \\
$F_1$ (weighted)     & ---             & 97,19\% \\
\bottomrule
\end{tabular}
\end{table}

Mô hình phân loại đúng 865 trên 890 ảnh kiểm thử, tức chỉ sai 25 ảnh.

Khoảng cách giữa kiểm định (98,19\%) và kiểm thử (97,19\%) là 1,00 điểm phần trăm.
Chênh lệch nhỏ và \emph{theo đúng chiều dự kiến} --- hiệu năng trên tập kiểm thử
thấp hơn --- là dấu hiệu lành mạnh: tập kiểm định đã bị dùng gián tiếp để chọn
điểm kiểm tra và điều chỉnh tốc độ học nên không còn khách quan, còn tập kiểm thử
chỉ được dùng đúng một lần. Nếu khoảng cách này lớn (chẳng hạn trên 5 điểm phần
trăm), đó sẽ là dấu hiệu mô hình đã bị điều chỉnh quá mức theo tập kiểm định.

\subsection{Ma trận nhầm lẫn}

\begin{figure}[H]
\centering
\includegraphics[width=0.8\textwidth]{figures/cls_confusion_matrix.png}
\caption[Ma trận nhầm lẫn trên tập kiểm thử]{Ma trận nhầm lẫn của mô hình
EfficientNet-B0 trên 890 ảnh kiểm thử. Phần lớn nhầm lẫn xảy ra giữa Algal Leaf
Spot và Phomopsis Leaf Spot.}
\label{fig:confusion}
\end{figure}

\begin{table}[H]
\centering
\caption{Ma trận nhầm lẫn trên tập kiểm thử (hàng: nhãn thật, cột: dự đoán)}
\label{tab:confusion}
\small
\begin{tabular}{l r r r r r | r}
\toprule
 & \textbf{Algal} & \textbf{Alloc.} & \textbf{Healthy} & \textbf{Blight} & \textbf{Phom.} & \textbf{Tổng} \\
\midrule
Algal Leaf Spot     & \textbf{142} & 1 & 1 & 0 & 3 & 147 \\
Allocaridara Attack & 2 & \textbf{180} & 0 & 0 & 1 & 183 \\
Healthy Leaf        & 1 & 0 & \textbf{195} & 0 & 0 & 196 \\
Leaf Blight         & 1 & 3 & 1 & \textbf{182} & 1 & 188 \\
Phomopsis Leaf Spot & 4 & 4 & 2 & 0 & \textbf{166} & 176 \\
\bottomrule
\end{tabular}
\end{table}

Ma trận nhầm lẫn xác nhận dự đoán đưa ra từ phân tích khám phá ở
\cref{sec:dulieu}: cặp \emph{Algal Leaf Spot} và \emph{Phomopsis Leaf Spot} nhầm
lẫn qua lại nhiều nhất --- 2 ảnh Algal bị nhận thành Phomopsis và 3 ảnh Phomopsis
bị nhận thành Algal, tổng cộng 5 trên 24 lỗi. Điều này phù hợp với quan sát hình
thái ở \cref{tab:hinhthai}: cả hai bệnh đều tạo đốm nhỏ trên nền lá xanh, khác
biệt chủ yếu ở kích thước và mật độ đốm.

Hai quan sát khác đáng lưu ý. \emph{Leaf Blight có precision tuyệt đối 100\%}:
cột Leaf Blight không chứa bất kỳ giá trị khác không nào nằm ngoài đường chéo ---
không một ảnh thuộc lớp khác từng bị nhận nhầm thành Leaf Blight. Mảng hoại tử lớn
màu nâu đen là dấu hiệu đặc thù mà không lớp nào khác có. Tuy vậy chính lớp này
lại có \emph{recall thấp nhất} (96,28\%): 7 ảnh Leaf Blight bị phân sang các lớp
khác, phân tán khá đều (2 sang Algal, 1 sang Allocaridara, 3 sang Healthy, 1 sang
Phomopsis). Cách phân bố lỗi này gợi ý rằng các ảnh Leaf Blight ở giai đoạn bệnh
sớm --- khi mảng hoại tử chưa lan rộng --- có biểu hiện gần với bệnh đốm hoặc thậm
chí gần lá khoẻ.

Đáng chú ý về mặt ứng dụng là hai chiều nhầm lẫn liên quan tới lớp Healthy Leaf.
Theo chiều \emph{báo động giả}, 5 ảnh lá khoẻ bị nhận thành lớp bệnh (1 sang Algal,
4 sang Phomopsis). Theo chiều \emph{bỏ sót}, 6 ảnh bệnh bị nhận thành lá khoẻ (2
từ Algal, 3 từ Leaf Blight, 1 từ Phomopsis).

Theo tiêu chí ưu tiên \en{recall} đã nêu ở \cref{sec:metrics}, chiều thứ nhất là
loại lỗi ít nghiêm trọng --- cảnh báo nhầm một cây khoẻ chỉ tốn thêm một lần kiểm
tra. Chiều thứ hai mới đáng lo, vì bỏ sót cây bệnh đồng nghĩa mất cơ hội can thiệp
sớm. Với 6 ảnh trên tổng số 694 ảnh bệnh trong tập kiểm thử, tỉ lệ bỏ sót bệnh là
0,86\% --- mức chấp nhận được cho ứng dụng thực tế.

\subsection{Hiệu năng theo từng lớp}

\begin{figure}[H]
\centering
\includegraphics[width=\textwidth]{figures/cls_per_class_metrics.png}
\caption[Chỉ số theo từng lớp]{Precision, recall và $F_1$ cho từng lớp trên tập
kiểm thử. Cả năm lớp đều đạt $F_1$ trên 95\%.}
\label{fig:per_class}
\end{figure}

\begin{table}[H]
\centering
\caption{Hiệu năng phân loại theo từng lớp trên tập kiểm thử}
\label{tab:cls_perclass}
\begin{tabular}{l c c c r}
\toprule
\textbf{Lớp} & \textbf{Precision} & \textbf{Recall} & \textbf{$F_1$} & \textbf{Số mẫu} \\
\midrule
Algal Leaf Spot     & 94,67\%  & 96,60\% & 95,62\% & 147 \\
Allocaridara Attack & 95,74\%  & 98,36\% & \textbf{97,04\%} & 183 \\
Healthy Leaf        & \textbf{97,99\%}  & \textbf{99,49\%} & \textbf{98,73\%} & 196 \\
Leaf Blight         & \textbf{100,00\%} & 96,81\% & 98,38\% & 188 \\
Phomopsis Leaf Spot & 97,08\%  & 94,32\% & 95,68\% & 176 \\
\midrule
\textbf{Trung bình có trọng số} & \textbf{97,22\%} & \textbf{97,19\%} & \textbf{97,19\%} & \textbf{890} \\
\bottomrule
\end{tabular}
\end{table}

Ba nhận xét từ \cref{tab:cls_perclass}:

\textbf{Không lớp nào bị bỏ rơi.} $F_1$ thấp nhất là 95,97\% (Algal Leaf Spot),
cao nhất 98,63\% (Allocaridara Attack) --- biên độ chỉ 2,66 điểm phần trăm. Kết
hợp với tỉ lệ mất cân bằng thấp (1,33), điều này cho thấy mô hình không thiên vị
lớp đa số.

\textbf{Recall vượt precision ở bốn trên năm lớp.} Ngoại lệ duy nhất là
Allocaridara Attack và Leaf Blight. Xu hướng chung này thuận lợi cho bài toán chẩn
đoán bệnh, theo tiêu chí đã nêu ở \cref{sec:metrics}.

\textbf{Algal Leaf Spot là lớp yếu nhất} với precision 94,70\% --- thấp nhất trong
năm lớp. Điều này nhất quán với ma trận nhầm lẫn và với phân tích màu ở
\cref{tab:meanrgb}, nơi lớp này có hiệu $G - R$ thấp nhất và phân bố kênh màu
chồng lấn nhiều nhất với các lớp còn lại.

\section{Phân tích bản đồ kích hoạt}
\label{sec:ketqua_xai}

Bản đồ kích hoạt được sinh cho toàn bộ \num{3104} ảnh huấn luyện, dùng điểm kiểm
tra tốt nhất từ \cref{sec:ketqua_cls}.

\begin{figure}[H]
\centering
\includegraphics[height=0.80\textheight, keepaspectratio]{figures/xai_gradcam_best.png}
\caption[Bản đồ kích hoạt Grad-CAM theo lớp]{Bản đồ kích hoạt của ảnh có độ tin
cậy cao nhất trong mỗi lớp. Từ trái sang phải: ảnh gốc, bản đồ nhiệt, và ảnh chồng
lớp.}
\label{fig:gradcam}
\end{figure}

\begin{table}[H]
\centering
\caption{Độ chính xác và độ tin cậy trên tập huấn luyện theo lớp}
\label{tab:gradcam_stats}
\begin{tabular}{l r c c c}
\toprule
\textbf{Lớp} & \textbf{Số ảnh} & \textbf{Accuracy} & \textbf{Tin cậy TB} & \textbf{Tin cậy nhỏ nhất} \\
\midrule
Algal Leaf Spot     & 513 & 99,03\% & 0,971 & \textbf{0,165} \\
Allocaridara Attack & 639 & 99,22\% & 0,981 & 0,220 \\
Healthy Leaf        & 683 & 99,56\% & 0,979 & 0,204 \\
Leaf Blight         & 655 & \textbf{99,85\%} & 0,988 & 0,354 \\
Phomopsis Leaf Spot & 614 & 99,84\% & \textbf{0,991} & 0,415 \\
\bottomrule
\end{tabular}
\end{table}

Độ chính xác trên tập huấn luyện đạt 99,03--99,85\%, cao hơn mức 97,19\% trên tập
kiểm thử. Chênh lệch này là bình thường và có hai nguyên nhân cộng hưởng: đây là
dữ liệu mô hình đã thấy khi huấn luyện, và khi suy luận thì không áp dụng tăng
cường dữ liệu nên bài toán dễ hơn.

Về mặt chất lượng bản đồ kích hoạt, quan sát \cref{fig:gradcam} cho thấy bốn kiểu
hành vi khác nhau:

\begin{itemize}
    \item \textbf{Algal và Phomopsis}: kích hoạt tập trung vào các đốm bệnh, chứng
          tỏ mô hình phân biệt hai lớp dựa trên đặc trưng đốm chứ không dựa vào
          màu tổng thể của lá.
    \item \textbf{Allocaridara}: kích hoạt bám vào vết cắn và lỗ thủng --- đây là
          lớp có bản đồ kích hoạt tập trung nhất, phù hợp với việc lớp này đạt
          $F_1$ cao nhất.
    \item \textbf{Leaf Blight}: kích hoạt trải rộng theo mảng hoại tử, không có
          tâm rõ ràng.
    \item \textbf{Healthy Leaf}: kích hoạt \emph{khuếch tán khắp phiến lá}, không
          có vùng trọng tâm.
\end{itemize}

Hành vi cuối cùng có ý nghĩa then chốt và cần được nhấn mạnh. Với lá khoẻ, mô hình
không có tổn thương nào để tập trung vào; nó nhận diện lớp này qua \emph{sự vắng
mặt} của dấu hiệu bệnh và qua màu xanh đồng đều trên toàn lá. Bản đồ kích hoạt vì
thế trải đều --- một hành vi hoàn toàn hợp lý xét theo nhiệm vụ phân loại, nhưng
lại vô nghĩa và gây hại nếu đem nhị phân hoá thành mặt nạ tổn thương. Đây chính là
bằng chứng thực nghiệm cho luận điểm lý thuyết ở \cref{sec:weaksup}: bản đồ kích
hoạt trả lời câu hỏi ``mô hình nhìn vào đâu'', không phải ``bệnh nằm ở đâu''.

Về độ tin cậy, giá trị trung bình đều trên 0,97 nhưng giá trị nhỏ nhất chênh lệch
đáng kể giữa các lớp: Algal Leaf Spot xuống tới 0,165 trong khi Phomopsis Leaf Spot
chỉ xuống 0,415. Điều này cho thấy Algal có một số trường hợp biên mà mô hình thực
sự không chắc chắn; bản đồ kích hoạt của chúng kém tin cậy hơn khi dùng làm nhãn
giả.

\section{Chất lượng nhãn giả}
\label{sec:ketqua_pseudo}

\subsection{Nhãn giả V1}

\begin{figure}[H]
\centering
\includegraphics[height=0.80\textheight, keepaspectratio]{figures/pl1_per_class.png}
\caption[Nhãn giả V1 theo lớp]{Mặt nạ giả V1 cho ảnh có độ tin cậy cao nhất mỗi
lớp. Từ trái sang phải: ảnh gốc, bản đồ nhiệt chồng lớp, mặt nạ giả chồng lớp.}
\label{fig:pl1}
\end{figure}

\begin{table}[H]
\centering
\caption{Thống kê độ phủ mặt nạ giả V1 theo lớp}
\label{tab:cov_v1}
\begin{tabular}{l c c c c}
\toprule
\textbf{Lớp} & \textbf{Trung bình} & \textbf{Độ lệch chuẩn} & \textbf{Nhỏ nhất} & \textbf{Lớn nhất} \\
\midrule
Algal Leaf Spot     & 15,1\% & 0,056 & 5,7\%  & 42,6\% \\
Allocaridara Attack & 13,7\% & 0,039 & 5,0\%  & 26,5\% \\
Healthy Leaf        & \textbf{26,4\%} & 0,076 & 11,1\% & \textbf{59,1\%} \\
Leaf Blight         & 12,3\% & 0,047 & 5,5\%  & 43,6\% \\
Phomopsis Leaf Spot & 14,5\% & 0,048 & 5,4\%  & 37,5\% \\
\bottomrule
\end{tabular}
\end{table}

Số liệu ở \cref{tab:cov_v1} định lượng chính xác vấn đề đã dự đoán ở
\cref{subsec:hanche_v1}. Lớp Healthy Leaf --- lớp \emph{không có} tổn thương ---
lại có độ phủ trung bình cao nhất (26,4\%), gần gấp đôi lớp Leaf Blight (12,3\%)
vốn là bệnh lan rộng nhất. Ảnh cực đoan nhất có tới 59,1\% diện tích bị đánh dấu
là ``bệnh'' trên một chiếc lá hoàn toàn khoẻ mạnh.

Độ lệch chuẩn cũng lớn (0,039--0,076), phản ánh việc ngưỡng cố định 0,5 không
thích ứng được với phân bố riêng của từng bản đồ nhiệt.

Toàn bộ \num{3104} mặt nạ được sinh trong 14 phút 00 giây. Độ chính xác dự đoán
trên tập này đạt 99,52\% với độ tin cậy trung bình 0,9822 --- nghĩa là chỉ 15 ảnh
có nhãn giả sinh ra từ dự đoán sai.

\subsection{Nhãn giả V2}

\begin{figure}[H]
\centering
\includegraphics[height=0.80\textheight, keepaspectratio]{figures/pl2_per_class.png}
\caption[Nhãn giả V2 theo lớp]{Mặt nạ giả V2 sinh bằng GradCAM++ đa tỉ lệ kết hợp
ngưỡng phân vị theo lớp. Mặt nạ của lớp Healthy Leaf là ma trận không.}
\label{fig:pl2}
\end{figure}

\begin{table}[H]
\centering
\caption{Độ phủ mục tiêu và độ phủ thực tế của nhãn giả V2}
\label{tab:cov_v2}
\begin{tabular}{l c c c c}
\toprule
\textbf{Lớp} & \textbf{Mục tiêu} & \textbf{Thực tế} & \textbf{Sai lệch} & \textbf{Độ lệch chuẩn} \\
\midrule
Algal Leaf Spot     & 18\% & 17,7\% & $-0{,}3$ & 0,008 \\
Allocaridara Attack & 20\% & 19,7\% & $-0{,}3$ & 0,007 \\
Healthy Leaf        & 0\%  & 0,0\%  & $0{,}0$  & 0,000 \\
Leaf Blight         & 22\% & 21,7\% & $-0{,}3$ & 0,007 \\
Phomopsis Leaf Spot & 18\% & 17,6\% & $-0{,}4$ & 0,009 \\
\bottomrule
\end{tabular}
\end{table}

Cơ chế ngưỡng phân vị hoạt động đúng như thiết kế: sai lệch so với mục tiêu không
vượt quá 0,4 điểm phần trăm ở bất kỳ lớp nào.

\begin{figure}[H]
\centering
\includegraphics[width=\textwidth]{figures/pl2_coverage.png}
\caption[Phân bố độ phủ V1 và V2]{Phân bố độ phủ của toàn bộ mặt nạ (trái) và độ
phủ trung bình theo lớp của V2 (phải). Cụm tại giá trị 0 trong biểu đồ bên trái
tương ứng 683 mặt nạ rỗng của lớp Healthy Leaf.}
\label{fig:pl2_coverage}
\end{figure}

\begin{table}[H]
\centering
\caption{So sánh tổng hợp nhãn giả V1 và V2}
\label{tab:v1v2_stats}
\begin{tabular}{l c c}
\toprule
\textbf{Chỉ số} & \textbf{V1} & \textbf{V2} \\
\midrule
Tổng số mặt nạ            & 3.104  & 3.104 \\
Độ chính xác dự đoán      & 99,52\% & 99,52\% \\
Độ tin cậy trung bình     & 0,9822 & 0,9822 \\
Độ phủ trung bình         & 0,1660 & 0,1504 \\
Mặt nạ gần rỗng ($<0{,}05$) & 1      & \textbf{683} \\
Mặt nạ quá lớn ($>0{,}60$)  & 0      & 0 \\
Thời gian sinh nhãn       & 14 ph 00 s & \textbf{4 ph 08 s} \\
\bottomrule
\end{tabular}
\end{table}

Hai chỉ số đầu giống hệt nhau giữa hai phiên bản, điều này là tất yếu vì cả hai
dùng chung một mô hình phân loại và cùng bộ ảnh --- chỉ khác cách chuyển bản đồ
kích hoạt thành mặt nạ.

Con số \textbf{683 mặt nạ gần rỗng} ở V2 cần được đọc đúng: đây \emph{không} phải
lỗi mà là kết quả có chủ đích của việc ép về không cho lớp Healthy Leaf theo
\cref{eq:hardzero}. Giá trị 683 trùng khớp chính xác với số ảnh lá khoẻ trong tập
huấn luyện, xác nhận cơ chế hoạt động đúng và không ảnh hưởng tới lớp nào khác.

Việc V2 nhanh hơn V1 tới 3,4 lần (4 phút 08 giây so với 14 phút 00 giây) dù phải
tính \emph{hai} bản đồ kích hoạt thay vì một là kết quả của cách cài đặt hiệu quả
hơn cho bước truyền ngược, không phải do giảm khối lượng tính toán.

\subsection{Đánh đổi đã được xác nhận}

Số liệu ở \cref{tab:cov_v2} xác nhận đánh đổi đã dự báo ở
\cref{subsec:tradeoff_v2}: độ lệch chuẩn của độ phủ giảm mạnh từ khoảng
0,039--0,076 ở V1 xuống chỉ còn 0,007--0,009 ở V2. Nói cách khác, \emph{mọi ảnh
trong cùng một lớp đều nhận mặt nạ có diện tích gần như bằng nhau}.

Đây vừa là thành công vừa là hạn chế. Thành công vì độ phủ được kiểm soát chặt chẽ
đúng như mục tiêu thiết kế. Hạn chế vì mặt nạ không còn mã hoá mức độ nặng nhẹ của
bệnh: một chiếc lá chỉ có một đốm nhỏ và một chiếc lá bị bệnh gần kín đều được gán
mặt nạ chiếm xấp xỉ 18\% hoặc 22\% diện tích. Mô hình phân đoạn huấn luyện trên
nhãn này học cách ``phủ đúng tỉ lệ diện tích'' thay vì học cách ``phát hiện đúng
vùng bệnh''.

\section{Phân tích tinh chỉnh bằng SAM}
\label{sec:ketqua_sam}

\subsection{Hiện tượng nở mặt nạ}

\begin{figure}[H]
\centering
\includegraphics[width=0.72\textwidth]{figures/sam_refinement.png}
\caption[So sánh nhãn V2 và nhãn SAM tinh chỉnh]{Bốn mẫu so sánh ảnh gốc, mặt nạ
giả V2 và mặt nạ sau khi SAM tinh chỉnh. SAM cho biên sắc nét hơn rõ rệt nhưng
đồng thời mở rộng diện tích mặt nạ đáng kể.}
\label{fig:sam_refine}
\end{figure}

Đo lường trực tiếp trên toàn bộ \num{3104} mặt nạ sau tinh chỉnh cho kết quả ở
\cref{tab:sam_expansion}.

\begin{table}[H]
\centering
\caption{Độ phủ trước và sau khi tinh chỉnh bằng SAM}
\label{tab:sam_expansion}
\begin{tabular}{l r r r r}
\toprule
\textbf{Lớp} & \textbf{Độ phủ V2} & \textbf{Độ phủ SAM} & \textbf{Độ lệch chuẩn} & \textbf{Mức tăng} \\
\midrule
Algal Leaf Spot     & 17,7\% & 34,63\% & 0,1210 & \textbf{+95,7\%} \\
Allocaridara Attack & 19,7\% & 31,98\% & 0,1112 & +62,3\% \\
Healthy Leaf        & 0,0\%  & 0,00\%  & 0,0000 & --- \\
Leaf Blight         & 21,7\% & 30,71\% & 0,1176 & +41,5\% \\
Phomopsis Leaf Spot & 17,6\% & 33,56\% & 0,1207 & +90,7\% \\
\midrule
\textbf{Toàn bộ}    & \textbf{15,04\%} & \textbf{25,43\%} & 0,1708 & \textbf{+69,1\%} \\
\bottomrule
\end{tabular}
\end{table}

Kết quả trái ngược với kỳ vọng ban đầu. \SAM{} được đưa vào quy trình với vai trò
\emph{tinh chỉnh biên} --- tức giữ nguyên vị trí và diện tích, chỉ làm biên chính
xác hơn. Thực tế nó \emph{mở rộng} diện tích mặt nạ từ 41,5\% đến 95,7\% tuỳ lớp,
trung bình toàn bộ tăng 69,1\%.

Hai chỉ số bổ sung làm rõ bản chất hiện tượng:

\begin{itemize}
    \item \textbf{200 trên \num{3104} mặt nạ (6,4\%) có độ phủ vượt 50\%} --- tức
          hơn một nửa diện tích ảnh bị đánh dấu là tổn thương, điều không thể xảy
          ra với bệnh đốm lá.
    \item \textbf{Độ lệch chuẩn tăng vọt} từ 0,007--0,009 (V2) lên 0,111--0,121
          (SAM), tức tăng khoảng 15 lần. Nhãn SAM biến thiên mạnh giữa các ảnh
          thay vì đồng nhất như V2.
    \item \textbf{683 mặt nạ rỗng được giữ nguyên rỗng} --- cơ chế ép về không cho
          lớp lá khoẻ không bị \SAM{} phá vỡ, đúng như thiết kế.
\end{itemize}

Nguyên nhân đã được dự báo về mặt lý thuyết ở \cref{sec:sam}: \SAM{} được huấn
luyện để phân đoạn \emph{đơn vị ngữ nghĩa tự nhiên} của cảnh. Khi nhận gợi ý là
vài điểm nằm trên một đốm bệnh, đơn vị ngữ nghĩa mà \SAM{} nhận diện không phải
đốm bệnh mà là \emph{cả chiếc lá} chứa đốm đó --- quan sát trực tiếp trên
\cref{fig:sam_refine} xác nhận điều này.

\subsection{Vì sao cơ chế bảo vệ IoU không phát hiện được}
\label{subsec:iouguard}

Cơ chế bảo vệ theo \cref{eq:iou_guard} lẽ ra phải phát hiện các mặt nạ sai lệch.
Nó đã không phát hiện được trường hợp nào của hiện tượng nở. Phân tích dưới đây
cho thấy điều này \emph{không phải sự cố ngẫu nhiên mà là hệ quả tất yếu} của công
thức được chọn.

\begin{quote}
\textbf{Mệnh đề.} Giả sử \SAM{} sinh ra mặt nạ $M_{\SAM}$ \emph{chứa trọn} mặt nạ
gốc $M_{\text{V2}}$, tức $M_{\text{V2}} \subseteq M_{\SAM}$ --- đúng với trường hợp
nở ra bao cả chiếc lá. Khi đó cơ chế bảo vệ với ngưỡng $\theta$ chỉ kích hoạt khi
độ phủ của \SAM{} vượt quá $1/\theta$ lần độ phủ của V2.
\end{quote}

\textbf{Chứng minh.} Khi $M_{\text{V2}} \subseteq M_{\SAM}$ thì
$|M_{\SAM} \cap M_{\text{V2}}| = |M_{\text{V2}}|$ và
$|M_{\SAM} \cup M_{\text{V2}}| = |M_{\SAM}|$. Do đó

\begin{equation}
    \operatorname{IoU}(M_{\SAM}, M_{\text{V2}})
    = \frac{|M_{\text{V2}}|}{|M_{\SAM}|}
    = \frac{c_{\text{V2}}}{c_{\SAM}}
    \label{eq:iou_containment}
\end{equation}

với $c$ ký hiệu độ phủ. Cơ chế bảo vệ kích hoạt khi
$c_{\text{V2}}/c_{\SAM} < \theta$, tương đương $c_{\SAM} > c_{\text{V2}}/\theta$.
$\square$

Áp dụng vào số liệu thực tế với $\theta = 0{,}15$ cho kết quả ở
\cref{tab:iouguard_analysis}.

\begin{table}[H]
\centering
\caption{Ngưỡng độ phủ cần thiết để cơ chế bảo vệ IoU kích hoạt}
\label{tab:iouguard_analysis}
\begin{tabular}{l c c c}
\toprule
\textbf{Lớp} & \textbf{Độ phủ V2} & \textbf{Độ phủ SAM cần vượt} & \textbf{Khả thi?} \\
\midrule
Algal Leaf Spot     & 17,7\% & 118,0\% & \textbf{Không} \\
Allocaridara Attack & 19,7\% & 131,3\% & \textbf{Không} \\
Leaf Blight         & 21,7\% & 144,7\% & \textbf{Không} \\
Phomopsis Leaf Spot & 17,6\% & 117,3\% & \textbf{Không} \\
\bottomrule
\end{tabular}
\end{table}

Kết luận rất rõ ràng: với mọi lớp bệnh, ngưỡng độ phủ cần thiết để kích hoạt cơ chế
bảo vệ đều \emph{vượt quá 100\% diện tích ảnh} --- một điều bất khả thi về mặt vật
lý. Nói cách khác:

\begin{quote}
Với cấu hình đã dùng, cơ chế bảo vệ dựa trên IoU \textbf{không bao giờ có thể}
phát hiện hiện tượng nở mặt nạ ở bất kỳ lớp bệnh nào. Nó chỉ chặn được trường hợp
\SAM{} \emph{trôi} sang một vùng hoàn toàn khác, chứ hoàn toàn bất lực trước trường
hợp \SAM{} \emph{nở} ra bao trùm mặt nạ gốc.
\end{quote}

Đây là một khiếm khuyết thiết kế mang tính hệ thống, không phải lỗi cài đặt. Nguồn
gốc của nó là IoU đối xứng --- nó phạt cả hai chiều sai lệch như nhau và do đó
không phân biệt được ``lệch vị trí'' với ``đúng vị trí nhưng quá rộng''.

\subsection{Giải pháp đề xuất}

Khiếm khuyết trên có thể khắc phục bằng cách bổ sung một ràng buộc \emph{bất đối
xứng} về tỉ lệ độ phủ, hoạt động song song với cơ chế IoU hiện có:

\begin{equation}
    M_{\text{refined}} =
    \begin{cases}
        M_{\text{V2}} & \text{nếu } \operatorname{IoU} < \theta \quad \text{(chống trôi)} \\[4pt]
        M_{\text{V2}} & \text{nếu } c_{\SAM} > \kappa \cdot c_{\text{V2}} \quad \text{(chống nở)} \\[4pt]
        M_{\SAM} & \text{trong các trường hợp còn lại}
    \end{cases}
    \label{eq:iou_guard_fixed}
\end{equation}

Với $\kappa = 2{,}0$, ràng buộc thứ hai sẽ kích hoạt cho phần lớn các trường hợp
nở đã quan sát --- chẳng hạn lớp Algal Leaf Spot có tỉ lệ nở trung bình 1,96 lần,
nằm sát ngưỡng, còn các ảnh cực đoan với độ phủ vượt 50\% đều bị chặn. Việc kiểm
chứng thực nghiệm cho giải pháp này nằm ngoài phạm vi luận văn và được đề xuất là
hướng phát triển ở \cref{chap:ketluan}.

\section{Kết quả phân đoạn}
\label{sec:ketqua_seg}

\subsection{Tổng hợp ba phiên bản}

\begin{table}[H]
\centering
\caption{Hiệu năng phân đoạn trên tập kiểm thử}
\label{tab:seg_results}
\small
\begin{tabular}{l c c c c c}
\toprule
\textbf{Phiên bản} & \textbf{IoU} & \textbf{Dice} & \textbf{Precision} & \textbf{Recall} & \textbf{Loss} \\
\midrule
V1 --- EfficientNet-UNet, nhãn V1$^{\dagger}$ & 0,7671 & 0,8623 & 0,8562 & 0,8877 & 0,3065 \\
V2 --- UNet++, nhãn V2                        & 0,5135 & 0,6301 & 0,6169 & 0,7419 & 0,5993 \\
V3 --- UNet++, nhãn SAM                       & 0,6239 & 0,7053 & 0,7090 & 0,8752 & 0,4704 \\
\bottomrule
\multicolumn{6}{l}{\small $^{\dagger}$ V1 dùng kiến trúc, phân bố nhãn và tập kiểm thử khác V2/V3 --- xem \cref{sec:caveat}.}
\end{tabular}
\end{table}

\begin{table}[H]
\centering
\caption{Diễn biến huấn luyện ba phiên bản phân đoạn}
\label{tab:seg_training}
\begin{tabular}{l C{1.7cm} C{2.5cm} C{2.5cm} C{2.5cm}}
\toprule
\textbf{Phiên bản} & \textbf{Số epoch} & \textbf{val\_loss nhỏ nhất} & \textbf{val\_IoU lớn nhất} & \textbf{Điểm kiểm tra chọn} \\
\midrule
V1 & 30 & 0,2959 (E23) & 0,7715 (E25) & E23 \\
V2 & 25 & 0,5957 (E20) & 0,5438 (E18) & E20 \\
V3 & 22 & 0,4478 (E12) & 0,6981 (E17) & E12 \\
\bottomrule
\end{tabular}
\end{table}

\begin{figure}[H]
\centering
\includegraphics[width=\textwidth]{figures/segv1_curves.png}
\caption[Đường cong huấn luyện phiên bản V1]{Diễn biến hàm mất mát, IoU và Dice
của phiên bản V1 qua 30 epoch.}
\label{fig:segv1_curves}
\end{figure}

\begin{figure}[H]
\centering
\includegraphics[width=\textwidth]{figures/segv2_curves.png}
\caption[Đường cong huấn luyện phiên bản V2]{Diễn biến huấn luyện của phiên bản V2
qua 25 epoch, kèm lịch tốc độ học cô-sin và biểu đồ so sánh với V1.}
\label{fig:segv2_curves}
\end{figure}

\begin{figure}[H]
\centering
\includegraphics[width=\textwidth]{figures/segv3_curves.png}
\caption[Đường cong huấn luyện phiên bản V3]{Diễn biến huấn luyện của phiên bản V3
qua 22 epoch, kèm biểu đồ so sánh ba phiên bản.}
\label{fig:segv3_curves}
\end{figure}

\subsection{Kết quả định tính}

\begin{figure}[H]
\centering
\includegraphics[width=\textwidth]{figures/segv1_predictions.png}
\caption[Dự đoán của phiên bản V1]{Bốn mẫu dự đoán của V1 trên tập kiểm thử: ảnh
đầu vào, nhãn giả tham chiếu, và mặt nạ dự đoán.}
\label{fig:segv1_pred}
\end{figure}

\begin{figure}[H]
\centering
\includegraphics[width=\textwidth]{figures/segv2_predictions.png}
\caption[Dự đoán của phiên bản V2]{Bốn mẫu dự đoán của V2: ảnh đầu vào, nhãn giả
V2, bản đồ xác suất và mặt nạ dự đoán kèm chỉ số IoU/Dice từng ảnh.}
\label{fig:segv2_pred}
\end{figure}

\begin{figure}[H]
\centering
\includegraphics[width=\textwidth]{figures/segv3_predictions.png}
\caption[Dự đoán của phiên bản V3]{Bốn mẫu dự đoán của V3 với nhãn tham chiếu là
mặt nạ đã được SAM tinh chỉnh. Biên dự đoán sắc nét hơn so với V2 nhưng diện tích
lớn hơn đáng kể.}
\label{fig:segv3_pred}
\end{figure}

\section{Phân tích tách bạch}
\label{sec:ablation}

\subsection{Giới hạn về khả năng so sánh giữa các phiên bản}
\label{sec:caveat}

Đọc \cref{tab:seg_results} một cách ngây thơ sẽ dẫn tới kết luận: V1 tốt nhất
(IoU 0,7671), V3 đứng thứ hai (0,6239), V2 kém nhất (0,5135) --- và do đó các
``cải tiến'' ở V2, V3 thực chất làm hỏng mô hình. \emph{Kết luận này sai}, và việc
chỉ ra vì sao nó sai là một trong những đóng góp chính của luận văn.

Vấn đề nằm ở chỗ mỗi phiên bản được đánh giá trên \emph{mốc tham chiếu của chính
nó}. Như đã nêu ở \cref{sec:metrics}, IoU và Dice luôn được định nghĩa tương đối
với một tập nhãn tham chiếu $Y$. Ở đây:

\begin{itemize}
    \item V1 được đo so với nhãn giả V1;
    \item V2 được đo so với nhãn giả V2;
    \item V3 được đo so với nhãn đã qua tinh chỉnh SAM.
\end{itemize}

Ba mốc tham chiếu này khác nhau về cả hình dạng, độ phủ lẫn độ biến thiên. So sánh
trực tiếp các con số IoU giữa chúng chẳng khác nào so sánh điểm số của ba học sinh
làm ba đề thi khác nhau với độ khó khác nhau.

\Cref{tab:confound} liệt kê đầy đủ các biến khác nhau giữa các cặp phiên bản.

\begin{table}[H]
\centering
\caption{Các biến khác biệt giữa các cặp phiên bản}
\label{tab:confound}
\small
\begin{tabular}{L{5.0cm} c c}
\toprule
\textbf{Biến} & \textbf{V1 so với V2} & \textbf{V2 so với V3} \\
\midrule
Phân bố nhãn tham chiếu   & Khác & Khác (chủ đích) \\
Kiến trúc mô hình         & Khác & Giống \\
Tập kiểm thử              & \textbf{Khác} & Giống \\
Hàm mất mát               & Khác & Giống \\
Bộ tối ưu và lịch tốc độ học & Khác & Giống \\
Bộ tăng cường dữ liệu     & Khác & Giống \\
Kích thước lô             & Khác & Khác (ngoài ý muốn) \\
Patience dừng sớm         & Khác & Khác (ngoài ý muốn) \\
\midrule
\textbf{Số biến khác nhau} & \textbf{8} & \textbf{3} \\
\bottomrule
\end{tabular}
\end{table}

Cặp V1--V2 khác nhau ở \emph{tám} biến cùng lúc. Với thiết kế thực nghiệm như vậy,
không thể quy chênh lệch hiệu năng cho bất kỳ nguyên nhân đơn lẻ nào. Đặc biệt,
việc hai phiên bản dùng \emph{hai tập kiểm thử khác nhau} --- do khác bộ sinh số
ngẫu nhiên như đã phân tích ở \cref{subsec:split_seg} --- khiến ngay cả phép so
sánh thô cũng mất cơ sở.

Vậy tại sao V1 lại đạt IoU cao nhất? Câu trả lời nằm ở \emph{độ khó của mục tiêu}
chứ không ở chất lượng mô hình:

\begin{itemize}
    \item Nhãn V1 là các khối tròn lớn, trơn, hình dạng đơn giản và dễ dự đoán ---
          hệ quả trực tiếp của việc phóng bản đồ $7 \times 7$ lên $224 \times 224$.
          Tái tạo lại một khối trơn như vậy là bài toán dễ.
    \item Nhãn V3 là các vùng do \SAM{} sinh ra: biên sắc, hình dạng phức tạp, bám
          theo cấu trúc thực của ảnh. Tái tạo chính xác các biên này khó hơn nhiều.
\end{itemize}

Nói cách khác, \textbf{IoU cao của V1 phản ánh việc mục tiêu của nó dễ, không phản
ánh việc nó định vị bệnh tốt hơn.} Trong điều kiện không có nhãn chuẩn ở mức điểm
ảnh, không thể thiết lập được hiệu năng định vị bệnh tuyệt đối của bất kỳ phiên
bản nào.

\subsection{Phép so sánh hợp lệ: V2 so với V3}

Cặp V2--V3 chỉ khác nhau ở ba biến, trong đó biến chính --- bộ nhãn huấn luyện ---
là biến ta muốn khảo sát. Kiến trúc, hàm mất mát, bộ tối ưu, bộ tăng cường và
\emph{tập kiểm thử} đều giống hệt nhau. Đây là phép so sánh có kiểm soát tốt nhất
trong luận văn.

\begin{table}[H]
\centering
\caption{So sánh V2 và V3 --- cùng kiến trúc, cùng tập kiểm thử}
\label{tab:v2v3}
\begin{tabular}{L{3.1cm} C{1.9cm} C{1.9cm} C{2.3cm} C{2.3cm}}
\toprule
\textbf{Chỉ số} & \textbf{V2} & \textbf{V3} & \textbf{Chênh tuyệt đối} & \textbf{Chênh tương đối} \\
\midrule
IoU       & 0,5135 & 0,6239 & $+0{,}1104$ & \textbf{+21,5\%} \\
Dice      & 0,6301 & 0,7053 & $+0{,}0752$ & +11,9\% \\
Precision & 0,6169 & 0,7090 & $+0{,}0921$ & +14,9\% \\
Recall    & 0,7419 & 0,8752 & $+0{,}1333$ & +18,0\% \\
\midrule
val\_loss nhỏ nhất & 0,5957 (E20) & 0,4478 (E12) & $-0{,}1479$ & $-24{,}8\%$ \\
\bottomrule
\end{tabular}
\end{table}

Cải thiện xuất hiện đồng đều trên cả bốn chỉ số, không có chỉ số nào đi ngược
chiều --- dấu hiệu của một tác động thực chứ không phải dao động ngẫu nhiên.

Bằng chứng bổ trợ mạnh mẽ đến từ tốc độ hội tụ: V3 đạt \texttt{val\_loss} 0,4478 ở
epoch 12, trong khi V2 phải tới epoch 20 mới đạt 0,5957 --- cao hơn đáng kể. Cùng
kiến trúc và cùng siêu tham số, mô hình học được nhanh hơn và sâu hơn trên nhãn
\SAM{}. Điều này cho thấy nhãn \SAM{} cung cấp \emph{tín hiệu huấn luyện sạch hơn}:
biên rõ ràng và nhất quán giúp mô hình tìm được quy luật dễ hơn.

\subsection{Ba yếu tố gây nhiễu cần trừ hao}
\label{subsec:confound}

Việc kết luận toàn bộ mức cải thiện +21,5\% là do \SAM{} sẽ là kết luận vội vàng.
Ba yếu tố sau có thể đóng góp một phần, và mục này đánh giá chiều tác động của
từng yếu tố.

\textbf{Yếu tố thứ nhất --- dừng sớm quá sớm ở V2.} Phiên bản V2 dùng patience = 5
còn V3 dùng patience = 10. Quan trọng hơn, dữ liệu cho thấy V2 \emph{vẫn đang cải
thiện} khi bị dừng: \texttt{val\_IoU} tăng từ 0,4960 tại epoch 20 lên 0,5339 tại
epoch 25 --- tức tăng 7,6\% trong 5 epoch cuối cùng trước khi dừng. Cơ chế dừng sớm
theo dõi \texttt{val\_loss} vốn đã chững lại, nên không nhận ra IoU vẫn đang lên.
Nếu V2 được cấp patience = 10 như V3, nhiều khả năng nó đã đạt hiệu năng cao hơn.
\emph{Chiều tác động: yếu tố này làm phóng đại khoảng cách V2--V3.}

\textbf{Yếu tố thứ hai --- kích thước mục tiêu.} Nhãn V3 có độ phủ trung bình
25,43\% so với 15,04\% của V2, tức lớn gấp 1,69 lần. Đây không phải chi tiết vô
hại, vì IoU có thiên lệch hệ thống theo kích thước đối tượng. Với một dải sai số
biên có bề rộng $w$ cố định, phần diện tích sai lệch tỉ lệ với chu vi nhân $w$,
trong khi mẫu số của IoU tỉ lệ với diện tích. Với các vùng có hình dạng tương tự
nhau, chu vi tỉ lệ với căn bậc hai của diện tích, nên sai số biên tương đối giảm
theo $1/\sqrt{A}$. Diện tích tăng 1,69 lần tương ứng sai số biên tương đối giảm
khoảng $1 - 1/\sqrt{1{,}69} \approx 23\%$. \emph{Chiều tác động: yếu tố này cũng
làm phóng đại lợi thế của V3.}

\textbf{Yếu tố thứ ba --- kích thước lô.} V2 dùng lô 4, V3 dùng lô 8. Với bộ nhớ
đồ hoạ chỉ 2,15 GB, đây là ràng buộc kỹ thuật chứ không phải lựa chọn thiết kế.
Lô lớn hơn cho ước lượng gradient ổn định hơn, có thể có lợi cho V3. \emph{Chiều
tác động: nhiều khả năng cũng có lợi cho V3, nhưng mức độ khó ước lượng.}

Cả ba yếu tố đều tác động \emph{cùng một chiều} là làm phóng đại lợi thế của V3.
Kết luận trung thực vì thế phải là:

\begin{quote}
Tinh chỉnh biên bằng \SAM{} \textbf{cải thiện chất lượng tín hiệu huấn luyện} ---
bằng chứng thuyết phục nhất là mô hình hội tụ nhanh hơn và đạt hàm mất mát kiểm
định thấp hơn hẳn với cùng kiến trúc và siêu tham số. Tuy nhiên, con số +21,5\%
IoU là \textbf{cận trên} của tác động này chứ không phải giá trị thực, vì cả ba
yếu tố gây nhiễu đã xác định đều tác động cùng chiều làm tăng khoảng cách.
\end{quote}

\subsection{Nghịch lý của SAM}

Một điểm thú vị về mặt phương pháp luận: \SAM{} vừa \emph{làm hỏng} nhãn theo tiêu
chí bệnh học --- mở rộng mặt nạ tới 69,1\%, bao cả những vùng lá khoẻ mạnh ---
lại vừa \emph{cải thiện} kết quả huấn luyện trên mọi chỉ số.

Hai sự thật này không mâu thuẫn. Chúng đo hai thứ khác nhau. Nhãn \SAM{} sai về
\emph{phạm vi} (bao quá rộng) nhưng đúng về \emph{cấu trúc} (biên bám theo ranh
giới thực trong ảnh, nhất quán giữa các ảnh). Mô hình học được từ tính nhất quán
cấu trúc đó, kể cả khi phạm vi bị phóng đại.

Hệ quả thực tiễn quan trọng: một mô hình phân đoạn huấn luyện trên nhãn V3 sẽ có
xu hướng dự đoán vùng bệnh rộng hơn thực tế. Với ứng dụng ước lượng \emph{tỉ lệ
diện tích lá bị bệnh} để ra quyết định canh tác, sai lệch hệ thống này phải được
hiệu chỉnh trước khi sử dụng.

\section{Phân tích lỗi và hạn chế}
\label{sec:phantichloi}

\subsection{Recall luôn vượt precision}

\begin{table}[H]
\centering
\caption{Quan hệ giữa recall và precision ở ba phiên bản}
\label{tab:recall_precision}
\begin{tabular}{l c c c}
\toprule
\textbf{Phiên bản} & \textbf{Precision} & \textbf{Recall} & \textbf{Chênh lệch} \\
\midrule
V1 & 0,8562 & 0,8877 & $+0{,}0315$ \\
V2 & 0,6169 & 0,7419 & $+0{,}1250$ \\
V3 & 0,7090 & 0,8752 & $+0{,}1662$ \\
\bottomrule
\end{tabular}
\end{table}

Cả ba phiên bản đều có recall vượt precision, và khoảng cách tăng dần từ V1 đến
V3. Điều này nghĩa là mô hình \emph{phân đoạn thừa} một cách hệ thống: nó đánh dấu
nhiều điểm ảnh là bệnh hơn so với nhãn tham chiếu.

Nguyên nhân là mô hình kế thừa đúng đặc tính của nhãn mà nó được huấn luyện. Nhãn
giả vốn đã có xu hướng bao rộng hơn tổn thương thực --- do bản đồ kích hoạt bị làm
mờ khi phóng từ độ phân giải thấp, và với V3 thì cộng thêm hiện tượng nở của
\SAM{}. Mô hình học đúng phân bố đó.

Xét theo yêu cầu ứng dụng nêu ở \cref{sec:metrics}, xu hướng này \emph{có lợi}:
trong chẩn đoán bệnh cây trồng, bỏ sót vùng bệnh nguy hiểm hơn cảnh báo thừa.

\subsection{Tiêu chí dừng sớm không phù hợp}
\label{subsec:earlystop}

Đây là hạn chế kỹ thuật nghiêm trọng nhất được phát hiện trong quá trình thực
nghiệm.

\begin{table}[H]
\centering
\caption{Diễn biến các epoch then chốt của phiên bản V3}
\label{tab:v3_epochs}
\begin{tabular}{c c c c l}
\toprule
\textbf{Epoch} & \textbf{val\_loss} & \textbf{val\_IoU} & \textbf{val\_Dice} & \textbf{Ghi chú} \\
\midrule
9  & 0,4514 & 0,6536 & 0,7348 & \\
11 & 0,4495 & 0,6711 & 0,7518 & \\
\textbf{12} & \textbf{0,4478} & 0,6479 & 0,7294 & \textbf{được chọn} (val\_loss nhỏ nhất) \\
16 & 0,4508 & 0,6917 & 0,7716 & \\
\textbf{17} & 0,4537 & \textbf{0,6981} & \textbf{0,7792} & \textbf{bị bỏ lỡ} (val\_IoU lớn nhất) \\
\bottomrule
\end{tabular}
\end{table}

Điểm kiểm tra được chọn theo tiêu chí \texttt{val\_loss} nhỏ nhất, rơi vào epoch
12 với val\_IoU = 0,6479. Nhưng epoch 17 có val\_IoU = 0,6981 --- \textbf{cao hơn
0,0502} --- lại bị bỏ qua vì val\_loss của nó cao hơn epoch 12 một lượng không
đáng kể (0,4537 so với 0,4478, chênh 1,3\%).

Nói cách khác, mô hình \emph{đã từng tốt hơn} trong quá trình huấn luyện nhưng
trạng thái đó không được lưu lại. Toàn bộ kết quả V3 báo cáo ở
\cref{tab:seg_results} vì thế là \emph{ước lượng thấp} so với khả năng thực của
phương pháp.

Nguyên nhân gốc đã được dự báo về mặt lý thuyết ở \cref{sec:loss}: hàm mất mát
FocalDice chủ động thu nhỏ gradient của các điểm ảnh dễ, khiến giá trị mất mát
tổng thể trở thành chỉ báo kém nhạy cho chất lượng phân đoạn. Hàm mất mát và IoU
không đồng pha với nhau.

Hiện tượng lặp lại ở V2 với mức độ còn nghiêm trọng hơn. Điểm kiểm tra được chọn
là epoch 20 với val\_IoU = 0,4960, trong khi epoch 18 đạt 0,5438. Đáng nói hơn,
epoch 20 là một \emph{điểm trũng cục bộ} của val\_IoU: các epoch lân cận đều cao
hơn hẳn (E18: 0,5438; E19: 0,5435; E21: 0,5379; E22: 0,5404). Nói cách khác, cơ
chế chọn theo val\_loss đã chọn đúng vào epoch có IoU kém nhất trong cả vùng.

\begin{table}[H]
\centering
\caption{Tổn thất do chọn điểm kiểm tra theo \texttt{val\_loss} ở ba phiên bản}
\label{tab:earlystop_loss}
\begin{tabular}{l l c c c c}
\toprule
\textbf{Phiên bản} & \textbf{Hàm mất mát} & \textbf{Epoch chọn} & \textbf{val\_IoU} & \textbf{val\_IoU tốt nhất} & \textbf{Tổn thất} \\
\midrule
V1 & BCE + Dice & E23 & 0,7699 & 0,7715 (E25) & $0{,}0016$ \\
V2 & FocalDice  & E20 & 0,4960 & 0,5438 (E18) & $\mathbf{0{,}0478}$ \\
V3 & FocalDice  & E12 & 0,6479 & 0,6981 (E17) & $\mathbf{0{,}0502}$ \\
\bottomrule
\end{tabular}
\end{table}

\Cref{tab:earlystop_loss} cho thấy một quy luật rất rõ và đây là bằng chứng thực
nghiệm trực tiếp cho dự báo lý thuyết ở \cref{sec:loss}. Phiên bản V1 dùng
BCE + Dice có tổn thất gần như bằng không (0,0016) --- hàm mất mát và IoU đồng pha
với nhau, chọn theo tiêu chí nào cũng cho kết quả tương đương. Hai phiên bản V2 và
V3 dùng FocalDice có tổn thất lớn gấp \emph{ba mươi lần} (0,0478 và 0,0502).

Cơ chế đằng sau đã được nêu ở \cref{sec:loss}: thành phần Focal chủ động thu nhỏ
gradient --- và qua đó thu nhỏ cả đóng góp vào giá trị hàm mất mát --- của những
điểm ảnh đã được phân loại tốt. Phần lớn diện tích ảnh là nền dễ, nên giá trị hàm
mất mát bị chi phối bởi một nhóm nhỏ điểm ảnh khó ở vùng biên. Trong khi đó IoU
tính trên toàn bộ vùng. Hai đại lượng vì thế đo hai khía cạnh khác nhau và không
còn đồng biến.

Khuyến nghị rút ra rất rõ ràng và có thể áp dụng ngay cho các nghiên cứu tương tự:
\emph{với các bài toán phân đoạn dùng hàm mất mát họ Focal, tiêu chí chọn điểm
kiểm tra và dừng sớm phải là chính chỉ số đánh giá mục tiêu} --- ở đây là val\_IoU
--- \emph{chứ không phải giá trị hàm mất mát}. Với hàm mất mát không có thành phần
điều biến theo độ khó như BCE + Dice, việc chọn theo hàm mất mát vẫn chấp nhận
được.

\subsection{Lỗi tăng cường dữ liệu ở phiên bản V1}

Trong cài đặt của V1, thuộc tính bật tăng cường được đặt trên đối tượng tập dữ
liệu \emph{dùng chung} cho cả ba tập con. Hệ quả là tập kiểm định và kiểm thử cũng
bị áp dụng tăng cường, dù chỉ gồm hai phép lật ngang và lật dọc.

Ảnh hưởng tới kết quả là nhỏ, vì phép lật bảo toàn diện tích và hình dạng của cả
ảnh lẫn mặt nạ, không làm thay đổi bản chất bài toán. Tuy nhiên đây vẫn là một sai
sót phương pháp cần ghi nhận: chỉ số của V1 mang thêm một thành phần dao động ngẫu
nhiên mà V2 và V3 không có. Đây là lý do thứ ba, bên cạnh khác mốc tham chiếu và
khác tập kiểm thử, khiến V1 không nên được dùng làm mốc so sánh.

\subsection{Dense CRF không khả dụng}

Như đã nêu ở \cref{sec:xulyanh}, Dense CRF là một hướng tinh chỉnh nhãn thay thế
cho \SAM{}, có ưu điểm là không cần mô hình nền tảng nặng và bám theo biên màu sắc
thực của ảnh.

Hướng này đã được thử nghiệm nhưng không đưa được vào quy trình: thư viện
\texttt{pydensecrf} không cài đặt được trên môi trường Python 3.13 dùng trong
nghiên cứu, do thư viện này chưa cung cấp bản dựng sẵn cho phiên bản Python đó.
Việc so sánh \SAM{} với Dense CRF trong vai trò bộ tinh chỉnh nhãn vì vậy vẫn là
một câu hỏi bỏ ngỏ, được đề xuất là hướng phát triển ở \cref{chap:ketluan}.

\subsection{Hạn chế của phép biến đổi đàn hồi}

Như đã lưu ý ở \cref{tab:augment}, bộ tăng cường của V2 và V3 gồm
\en{ElasticTransform} với xác suất 0,3 và \en{GridDistortion} với xác suất 0,2.
Hai phép này làm biến dạng \emph{đồng thời} cả ảnh lẫn mặt nạ.

Với nhãn chuẩn, đây là cách tăng cường hợp lệ. Với nhãn giả vốn đã không chắc chắn
về biên, việc biến dạng thêm có thể làm nhiễu tín hiệu huấn luyện thay vì tăng khả
năng khái quát. Giả thuyết này chưa được kiểm chứng bằng thực nghiệm tách bạch
trong luận văn, và được ghi nhận là một hướng khảo sát.

\section{Kết luận chương}

Chương này đã trình bày và phân tích toàn bộ kết quả thực nghiệm. Các phát hiện
chính:

\begin{enumerate}
    \item \textbf{Phân loại đạt 97,19\% trên tập kiểm thử} với $F_1$ 97,19\%, cả
          năm lớp đều vượt $F_1$ 95\%. Cặp nhầm lẫn chính là Algal Leaf Spot và
          Phomopsis Leaf Spot, đúng như dự đoán từ phân tích khám phá.

    \item \textbf{Bản đồ kích hoạt của lớp lá khoẻ khuếch tán khắp phiến lá}, xác
          nhận bằng thực nghiệm rằng bản đồ kích hoạt phản ánh cơ sở ra quyết định
          của mô hình chứ không phải vị trí tổn thương --- cơ sở cho biện pháp ép
          mặt nạ lớp này về không.

    \item \textbf{Ngưỡng phân vị đạt độ phủ mục tiêu với sai lệch dưới 0,4 điểm
          phần trăm}, đồng thời làm độ lệch chuẩn của độ phủ giảm khoảng năm lần
          --- xác nhận cả thành công lẫn đánh đổi đã dự báo.

    \item \textbf{SAM mở rộng độ phủ mặt nạ 69,1\%} thay vì chỉ tinh chỉnh biên,
          với 200 mặt nạ vượt 50\% diện tích ảnh.

    \item \textbf{Cơ chế bảo vệ dựa trên IoU không thể phát hiện hiện tượng nở} ---
          đã chứng minh bằng lập luận toán học rằng ngưỡng cần thiết để kích hoạt
          vượt quá 100\% diện tích ảnh với mọi lớp bệnh. Một ràng buộc bất đối xứng
          về tỉ lệ độ phủ được đề xuất để khắc phục.

    \item \textbf{Tinh chỉnh bằng SAM cải thiện V2 lên V3 trên mọi chỉ số}
          (IoU +21,5\%), nhưng ba yếu tố gây nhiễu đều tác động cùng chiều nên con
          số này là cận trên chứ không phải giá trị thực.

    \item \textbf{Chỉ số của V1 không so sánh được với V2 và V3} do khác biệt ở
          tám biến, bao gồm cả tập kiểm thử. IoU cao của V1 phản ánh mục tiêu dễ,
          không phản ánh chất lượng định vị bệnh.

    \item \textbf{Tiêu chí dừng sớm theo val\_loss là không phù hợp} với hàm mất
          mát họ Focal --- V3 bỏ lỡ điểm kiểm tra tốt hơn 0,0502 IoU.
\end{enumerate}

\Cref{chap:ketluan} tổng hợp các kết quả này thành câu trả lời trực tiếp cho bốn
câu hỏi nghiên cứu đặt ra ở \cref{sec:muctieu}, đồng thời đề xuất hướng phát triển.
